% Stammfunktion und Hauptsatz – bank/stammfunktion-und-hauptsatz/ (zusammenbau v0.4, Heft)
\einheitenkopf[t1]{Stammfunktion und Hauptsatz}
\setcounter{aufgabe}{0}

% Anlauf (ohne Prüfkennung)
\begin{aufgabe}{Stammfunktion – Anlauf}
\begin{teile}
\teil Gesucht ist eine Funktion $F$ mit $F'(x) = 2x^3 - 5$. Ableiten oder aufleiten? \\ \kreuz{ableiten} \\ \kreuz{aufleiten}
\teil $f(x) = x \cdot (x + 4)$ – eine Stammfunktion $F$? $F(x) = $ \leerfeld
\teil $F(x) = x^2 - 2x$ ist eine Stammfunktion von $f$. Gib alle Stammfunktionen von $f$ an, deren Graph überall den senkrechten Abstand $3$ zum Graphen von $F$ hat.
\end{teile}
\end{aufgabe}

\begin{aufgabe}{Stammfunktion – Prüfungsaufgaben}
\begin{teile}
\teil $F(x) = 3x^4 - 2x + 5$ ist eine Stammfunktion von $f$. Gib $f(x)$ an. \hfill (Abitur 2020 GK) \\ $f(x) = $ \leerfeld
\teil $F(x) = \frac{1}{5}x^5 + x^2 - 7$ ist eine Stammfunktion von $f$. Gib $f(x)$ an. \hfill (Abitur 2020 GK) \\ $f(x) = $ \leerfeld
\teil $F(x) = 2e^{3x} + x$ ist eine Stammfunktion von $f$. Gib $f(x)$ an. \hfill (Abitur 2020 GK) \\ $f(x) = $ \leerfeld
\teil Weise nach, dass $F(x) = (-x - 3) \cdot e^{-x}$ eine Stammfunktion von $f(x) = (x + 2) \cdot e^{-x}$ ist, und begründe, dass $f$ nur bei $x = -2$ eine Nullstelle hat. \hfill (Abitur 2018 GK)
\teil Weise nach, dass $F(x) = (x^2 - 4x + 4) \cdot e^{x}$ eine Stammfunktion von $f(x) = x \cdot (x - 2) \cdot e^{x}$ ist. \hfill (Abitur 2020 GK)
\teil $f(x) = 6x^2 - 4x$. Bestimme die Stammfunktion $F$ von $f$, deren Graph durch $P(2 | 5)$ verläuft. \hfill (Abitur 2025 GK) \\ $F(x) = $ \leerfeld
\teil $f(x) = \frac{1}{2}x^3 - 1$. Der Punkt $Q(-2 | 3)$ liegt auf dem Graphen einer Stammfunktion $F$ von $f$. Bestimme $F(x)$. \hfill (Abitur 2026 GK) \\ $F(x) = $ \leerfeld
\teil $F(x) = (2x + 3) \cdot e^{-x}$ ist eine Stammfunktion von $f$. Bestimme die Stammfunktion $H$ von $f$ mit $H(0) = 8$. \hfill (Abitur 2020 GK) \\ $H(x) = $ \leerfeld
\teil $F(x) = x^3 - 3x$ ist eine Stammfunktion von $f(x) = 3x^2 - 3$. Bestimme die beiden Stammfunktionen von $f$, deren Graph die x-Achse berührt. \hfill (Abitur 2024 GK)
\teil $f(x) = (x^2 - 3) \cdot e^{x}$. Es gilt: Das Integral über $x^2 \cdot e^{x}$ ist $(x^2 - 2x + 2) \cdot e^{x} + C$. Gib damit eine Stammfunktion von $f$ an. \hfill (Abitur 2023 GK) \\ $F(x) = $ \leerfeld
\teil $f(x) = (3x - 1) \cdot e^{-x}$. Es gilt: Das Integral über $x \cdot e^{-x}$ ist $-(x + 1) \cdot e^{-x} + C$. Gib damit eine Stammfunktion von $f$ an. \hfill (Abitur 2023 GK) \\ $F(x) = $ \leerfeld
\end{teile}
\end{aufgabe}

% Anlauf (ohne Prüfkennung)
\begin{aufgabe}{Hauptsatz – Anlauf}
\begin{teile}
\teil Berechne das Integral von $0$ bis $2$ über $f(x) = 3x^2 - 2x$. Wert: \leerfeld
\teil $F(x) = \sqrt{2x + 1}$ ist eine Stammfunktion von $f$. Berechne das Integral von $0$ bis $4$ über $f(x)$. Wert: \leerfeld
\end{teile}
\end{aufgabe}

\begin{aufgabe}{Hauptsatz – Prüfungsaufgaben}
\begin{teile}
\teil Berechne das Integral von $0$ bis $2$ über $f(x) = x^3 - 3x^2$. \hfill (Abitur 2023 GK) \\ Wert: \leerfeld
\teil Berechne das Integral von $0$ bis $2$ über $f(x) = \frac{1}{4} \cdot (x^3 - 6x^2 + 8x)$. \hfill (Abitur 2018 GK) \\ Wert: \leerfeld
\teil Berechne das Integral von $0$ bis $2\pi$ über $f(x) = \mathrm{cos}(x) + 2$. \hfill (Abitur 2021 GK) \\ Wert: \leerfeld
\teil Berechne das Integral von $0$ bis $2\pi$ über $f(x) = 3 - 2 \cdot \mathrm{sin}(x)$. \hfill (Abitur 2021 GK) \\ Wert: \leerfeld
\teil Berechne das Integral von $0$ bis $\pi$ über $f(x) = \mathrm{sin}(2x) + 1$. \hfill (Abitur 2021 GK) \\ Wert: \leerfeld
\teil $F(x) = (2 - x) \cdot e^{x}$ ist eine Stammfunktion von $f(x) = (1 - x) \cdot e^{x}$. Ein Näherungswert für das Integral von $-2$ bis $1$ über $f(x)$ ist $2{,}1$. Berechne den exakten Wert und die prozentuale Abweichung des Näherungswerts. \hfill (Abitur 2025 GK) \\ exakt: \leerfeld , Abweichung: \leerfeld \%
\teil $F(x) = (x - 2) \cdot e^{\frac{x}{2}}$ ist eine Stammfunktion von $f(x) = \frac{x}{2} \cdot e^{\frac{x}{2}}$. Ein Näherungswert für das Integral von $0$ bis $2$ über $f(x)$ ist die Dreiecksfläche $\frac{1}{2} \cdot 2 \cdot f(2) = e$. Berechne den exakten Wert und die prozentuale Abweichung des Näherungswerts. \hfill (Abitur 2025 GK) \\ exakt: \leerfeld , Abweichung: \leerfeld \%
\teil $f(x) = 3x - \frac{1}{4}x^3$. Ein Näherungswert für das Integral von $0$ bis $2$ über $f(x)$ ist die Dreiecksfläche $\frac{1}{2} \cdot 2 \cdot f(2) = 4$. Berechne die prozentuale Abweichung des Näherungswerts vom exakten Wert. \hfill (Abitur 2022 GK) \\ Abweichung: \leerfeld \%
\teil $f(x) = 4x - x^2$. Ein Näherungswert für das Integral von $0$ bis $3$ über $f(x)$ ist die Rechteckfläche $3 \cdot f(1{,}5)$. Berechne die prozentuale Abweichung des Näherungswerts vom exakten Wert. \hfill (Abitur 2022 GK) \\ Abweichung: \leerfeld \%
\end{teile}
\end{aufgabe}

% Anlauf (ohne Prüfkennung)
\begin{aufgabe}{Der Graphenblick – Anlauf}
\begin{teile}
\teil $u(x) = x^3 + 2$ und $v(x) = 3x^2$. Eine der beiden ist Stammfunktion der anderen. Welche ist die Stammfunktion? \\ \kreuz{$u$} \\ \kreuz{$v$}
\teil Das Bild zeigt den Graphen von $f$. Skizziere den Graphen der Stammfunktion $F$ von $f$, die durch $P(0 | 1)$ geht.

\begin{ksys}[xmin=-2,xmax=4,ymin=-3,ymax=4] \gerade{1}{-1}{f} \punkt{0}{1}{P} \end{ksys}
\teil Das Bild zeigt den Graphen einer Stammfunktion $F$ von $f$. Lies passende Werte ab und berechne das Integral von $-1$ bis $3$ über $f(x)$.

\begin{ksys}[xmin=-2,xmax=4,ymin=-1,ymax=5,ablesen] \parabel{-0.5}{1}{4}{F} \end{ksys}
Wert: \leerfeld
\end{teile}
\end{aufgabe}

\begin{aufgabe}{Der Graphenblick – Prüfungsaufgaben}
\begin{teile}
\teil Das Bild zeigt den Graphen von $f$; es gilt $f(x) > 0$, der Hochpunkt ist $(0 | 2)$. Skizziere den Graphen der Stammfunktion $F$ von $f$, die durch $P(0 | 1)$ geht. \hfill (Abitur 2023 GK)

\begin{ksys}[xmin=-4,xmax=4,ymin=-2,ymax=4] \funktion{2/(1+\x^2)}{f} \punkt{0}{1}{P} \end{ksys}
\teil Das Bild zeigt den Graphen von $f$. Skizziere nur mithilfe dieses Graphen den Graphen der Stammfunktion $F$ von $f$, die durch $P(1 | -1)$ geht. \hfill (Abitur 2021 GK)

\begin{ksys}[xmin=-4,xmax=3,ymin=-2,ymax=4] \funktionab{(\x-1)*exp(\x-1)}{f}{-4}{2.2} \punkt{1}{-1}{P} \end{ksys}
\teil Das Bild zeigt den Graphen einer Stammfunktion $F$ von $f$. Lies passende Werte ab und berechne das Integral von $0$ bis $4$ über $f(x)$. \hfill (Abitur 2021 GK) \\

\begin{ksys}[xmin=-1,xmax=6,ymin=-1,ymax=5,ablesen] \parabel{0.5}{3}{-0.5}{F} \end{ksys}
Wert: \leerfeld
\teil Begründe, dass jede Stammfunktion von $g(x) = x^2 \cdot (x - 2)$ bei $x = 2$ ihr Minimum annimmt. \hfill (Abitur 2026 GK)
\teil $h(x) = -\frac{1}{2}x \cdot (x - 6)$. An welcher Stelle mit $x > 0$ nimmt jede Stammfunktion von $h$ ihr Maximum an? Begründe ohne Stammfunktionsterm. \hfill (Abitur 2019 GK) \\ $x = $ \leerfeld
\teil Das Bild zeigt die Graphen von $f$ und $f'$. Aussage: Der Graph jeder Stammfunktion von $f$ ändert genau einmal sein Krümmungsverhalten. Beurteile die Aussage. \hfill (Abitur 2024 GK) \\ \kreuz{wahr} \\ \kreuz{falsch}

\begin{ksys}[xmin=-3,xmax=3,ymin=-3,ymax=6,ablesen] \funktionab{\x^3/3-\x}{f}{-2.6}{2.6} \parabel{1}{0}{-1}{f'} \end{ksys}
\end{teile}
\end{aufgabe}

\begin{aufgabe}{Der Graphenblick – Prüfungsaufgaben – weiter}
\begin{teile}
\teil $f(x) = x \cdot (x + 2) \cdot e^{-x}$; $F(x) = -(x + 2)^2 \cdot e^{-x}$ ist eine Stammfunktion von $f$. Begründe ohne den Term von $F$, dass der Graph jeder Stammfunktion von $f$ einen Tiefpunkt auf der y-Achse hat, und bestimme die Stammfunktion $H$, deren Tiefpunkt im Ursprung liegt. \hfill (Abitur 2020 GK) \\ $H(x) = $ \leerfeld
\teil $f(x) = \frac{1}{2}x \cdot (4 - x) \cdot e^{x}$; $F(x) = \frac{1}{2} \cdot (-x^2 + 6x - 6) \cdot e^{x}$ ist eine Stammfunktion von $f$. Begründe ohne den Term von $F$, dass der Graph jeder Stammfunktion von $f$ einen Tiefpunkt auf der y-Achse hat, und bestimme die Stammfunktion $H$, deren Tiefpunkt im Ursprung liegt. \hfill (Abitur 2020 GK) \\ $H(x) = $ \leerfeld
\end{teile}
\end{aufgabe}
